% !TeX root = ../papers/nc-p3.tex
\section{孤立した両側直線ファイバーの一次変形による排除}
\label{sec:exclude-two-sided-centers}
\begin{theorem}\label{thm:no-two-sided-centers}
標準正則窓で前点も後続点も直線をなす点は存在しない。
従って直線ファイバーを持つ点には逆側にただ一つの点があり、その逆側評価階数は3である。
\end{theorem}
\begin{proof}[共通の階数二の一次変形]
前節の両側点 $p$ を仮定し、$X,F,S,T,a,H$ を同じ記号とする。
$p$ を含む射影平面 $\mathbb P(H)$ 内での一次近傍を $P$ とする。その台は一点、
局所環は $k[\xi,\eta]/(\xi,\eta)^2$ である。$p$ の任意のこの一次変形について、
順・逆支持は同時に階数2の直和因子に持ち上がる。

実際 $v\in H$ について $w(v)$ の $X,F$ 両側への射影は $v(a)h=0$ である。
$X$ 側だけへの射影は $X\otimes(R_2\cap(V_2\otimes T))$ に入る。
括弧内は $w(p)$ の位置1からの縮約で張られる二空間なので、この射影は
一意な $B_v:S\to X$ を $w(p)$ の位置1に作用させたものである。
同様に $F$ 側だけの射影には $C_v:T\to F$ がある。補空間を固定して持ち上げれば、
$p+\epsilon v$ での順支持は $s\mapsto s+\epsilon B_v(s)$ のグラフ、逆支持も
$t\mapsto t+\epsilon C_v(t)$ のグラフである。
包含は表示式から、等号は閉点での階数と中山の補題から従う。
二方向の一次パラメータの全積が零なので、これらは $P$ 上の支持部分束 $S_P,T_P$ を与える。

商束を $E=V_1/S_P,F'=V_3/T_P$ とする。
$R_1\to E\otimes V_2$, $R_2\to V_2\otimes F'$ の像 $I,J$ は階数4の直和因子である。
例えば第一の核には、変形した $w$ を最後から縮約した二つの独立元があり、
閉点でこれらは核全体を張る。階数4の可逆小行列とこの核を使うと写像を
四次の単位行列と零の和に消去できる。第二も同じである。
\end{proof}

\begin{proof}[始点と終点を固定した十二文字]
$P\times P$ 上の独立な両境界について、三文字挿入の正方行列
\[
 (I\boxtimes F')\oplus(E\boxtimes J)\longrightarrow E\boxtimes F'\otimes V_2
\]
を取る。前節により閉点で階数15である。可逆小行列による消去で、その余核は一つの関数 $f$
の剰余環として表示でき、$f(p,p)=0$ である。特に双対数上で
$(p+\epsilon u,p+\epsilon v)$ に沿って可逆な一文字分の商が得られる条件は
$L(u)+M(v)=0$ の一一次式である。$L,M$ は二次元接空間上の線型形式で、零でもよい。

今、連続する四つの境界で $p, p+\epsilon u,p+\epsilon v,p$ と置く。
三つの三文字商を全て可逆とする条件は
\[
           M(u)=0,\qquad L(u)+M(v)=0,\qquad L(v)=0 .
\]
未知数四個に対する三本の斉次式なので $(u,v)\ne(0,0)$ を取れる。
各三文字では可逆な余核への像が単元となる一元を選ぶ。
境界の一次式をその点で評価し、これらを順に掛けると、
両端を元の $p$ に固定した経路の評価が得られる。内部関係は余核の定義で、境界関係は
$S_P,T_P$ の定義で零である。非零方向のある中間境界では二つの文字の評価が
$k[\epsilon]/\epsilon^2$ の全体を張る（定数と非零一次項を独立に取れる）。
他の境界では単元評価を固定すればよい。

この十二文字の区間を同じ固定点でつなぎ、各区間に独立な双対数 $\epsilon_j$ を与える。
元の窓の $12n+1$ 文字を $\bigotimes_{j=1}^n k[\epsilon_j]/\epsilon_j^2$ に評価する写像は
全関係を消し、上の二つの選択の積により全射である。従って次元は $2^n$ 以上となり、
標準Hilbert関数に反する。逆側が空でないことは命題\ref{prop:isolated-branch-blocks}による。
\end{proof}
